%% POUR LE CORRIGÉ : remplacer \corrigefalse par \corrigetrue (dans \begin{document}).
%% Avec \corrigefalse, on obtient le sujet seul (1 page) ; avec \corrigetrue,
%% la même page avec les réponses en rouge (et les côtés repassés en rouge).

%% Font size %%
\documentclass[11pt]{article}

%% Load the custom package
\usepackage{Mathdoc}
\usepackage{colortbl}
\usetikzlibrary{calc}

%% Numéro de séquence %% Titre de la séquence %%
\renewcommand{\centerhead}{S04 - Éval : Sommet et côtés d'un angle}

%% Spacing commands %%
\renewcommand{\baselinestretch}{1} \setlength{\parindent}{0pt}

% Corrigé : la feuille est écrite une seule fois (\feuille).
\newif\ifcorrige
% Réponse (en mode math) dans une case de largeur fixe : \rep[largeur]{réponse}
% (pointillés sur le sujet, réponse en rouge sur le corrigé, même place)
\newcommand{\rep}[2][0.8cm]{\makebox[#1]{\ifcorrige\textcolor{red}{$#2$}\else\dotfill\fi}}
% Écriture à entourer : \entoure{faux} ou \entoure[r]{bonne réponse}
\newcommand{\entoure}[2][]{\def\coulcadre{white}\ifcorrige\if r#1\def\coulcadre{red}\fi\fi\tikz[baseline=(m.base)]\node[draw=\coulcadre,thick,rounded corners=2mm,inner sep=2pt](m){#2};}
% Écriture entourée en noir (exemple déjà fait)
\newcommand{\entoureex}[1]{\tikz[baseline=(m.base)]\node[draw=black,thick,rounded corners=2mm,inner sep=2pt](m){#1};}
% Notation d'un angle
\newcommand{\ang}[1]{\widehat{#1}}

% Triangle XYZ dont les côtés sont prolongés (droites fines) ;
% dans le corrigé, les côtés de l'angle de sommet #1 sont repassés en rouge.
% #7, #8, #9 : direction (en degrés) de chaque étiquette par rapport à son sommet.
% \triangleCotes{X}{Y}{Z}{x_X,y_X}{x_Y,y_Y}{x_Z,y_Z}{dir_X}{dir_Y}{dir_Z}
\newcommand{\triangleCotes}[9]{%
  \begin{tikzpicture}[scale=0.62]
    \coordinate (P1) at (#4); \coordinate (P2) at (#5); \coordinate (P3) at (#6);
    \draw[gray] ($(P1)!-0.35!(P2)$) -- ($(P2)!-0.35!(P1)$);
    \draw[gray] ($(P2)!-0.35!(P3)$) -- ($(P3)!-0.35!(P2)$);
    \draw[gray] ($(P3)!-0.35!(P1)$) -- ($(P1)!-0.35!(P3)$);
    \draw[thick] (P1) -- (P2) -- (P3) -- cycle;
    \ifcorrige
      \draw[red, line width=1.6pt] ($(P2)!-0.35!(P1)$) -- (P1) -- ($(P3)!-0.35!(P1)$);
    \fi
    \node at ($(P1)+(#7:4.5mm)$) {$#1$};
    \node at ($(P2)+(#8:4.5mm)$) {$#2$};
    \node at ($(P3)+(#9:4.5mm)$) {$#3$};
  \end{tikzpicture}%
}

% Trois demi-droites de même origine T passant par F, G, D ;
% dans le corrigé, les côtés de l'angle GTD sont repassés en rouge.
\newcommand{\eventail}{%
  \begin{tikzpicture}[scale=0.62]
    \coordinate (T) at (0,0);
    \coordinate (F) at (10:2.2); \coordinate (G) at (85:1.9); \coordinate (D) at (150:2.0);
    \draw[thick] (T) -- (10:3.5); \draw[thick] (T) -- (85:3.0); \draw[thick] (T) -- (150:3.1);
    \ifcorrige
      \draw[red, line width=1.6pt] (85:3.0) -- (T) -- (150:3.1);
    \fi
    \foreach \p in {F,G,D} \draw (\p) ++(-2pt,-2pt) -- ++(4pt,4pt) ++(-4pt,0) -- ++(4pt,-4pt);
    \node at ($(F)+(-80:4mm)$) {$F$};
    \node at ($(G)+(-5:4mm)$) {$G$};
    \node at ($(D)+(240:4mm)$) {$D$};
    \node at (-90:3.5mm) {$T$};
  \end{tikzpicture}%
}

\newcommand{\feuille}{%
\phantom{0}
\vspace{-.5cm}

\entetedevoirs{20}

\begin{center}
\duree{15 minutes}
\total{20 points}
\calculatrice{0}
\end{center}

\vspace{-0.7cm}
\begin{exercicedevoir}[3][Trouve le sommet]
\vspace{-0.3cm}
Écris le \textbf{sommet} de chaque angle. \hfill \textit{\small (0,5 point par réponse)}

\renewcommand{\arraystretch}{1.5}
\begin{tabularx}{\textwidth}{@{}XXXX@{}}
\textit{Ex.} $\ang{KLM}$ : \makebox[0.8cm]{$L$} &
a) $\ang{PQR}$ : \rep{Q} &
b) $\ang{HIJ}$ : \rep{I} &
c) $\ang{CAB}$ : \rep{A} \\
 & d) $\ang{YZW}$ : \rep{Z} &
e) $\ang{NMV}$ : \rep{M} &
f) $\ang{XEU}$ : \rep{E} \\
\end{tabularx}
\end{exercicedevoir}

\vspace{-1.1cm}
\begin{exercicedevoir}[5][Complète le tableau]
\vspace{-0.3cm}
\begin{minipage}[c]{0.3\linewidth}
Complète le tableau.

\smallskip
\textbf{Attention :} sur les lignes d et e, c'est le \textbf{nom de l'angle} qu'il faut trouver !

\smallskip
\textit{\small (1 point par ligne)}
\end{minipage}\hfill
\begin{minipage}[c]{0.68\linewidth}
\centering
\renewcommand{\arraystretch}{1.3}
\begin{tabular}{|l|c|c|}
\hline
\multicolumn{1}{|c|}{\textbf{Angle}} & \textbf{Sommet} & \textbf{Côtés} \\ \hline
\rowcolor{gray!15} \textit{Ex.} \ $\ang{XYZ}$ & $Y$ & $[YX)$ \ et \ $[YZ)$ \\ \hline
a) \ $\ang{BCD}$ & \rep{C} & \rep[1.3cm]{[CB)} \ et \ \rep[1.3cm]{[CD)} \\ \hline
b) \ $\ang{STU}$ & \rep{T} & \rep[1.3cm]{[TS)} \ et \ \rep[1.3cm]{[TU)} \\ \hline
c) \ $\ang{FHG}$ & \rep{H} & \rep[1.3cm]{[HF)} \ et \ \rep[1.3cm]{[HG)} \\ \hline
d) \ \rep[3.3cm]{\ang{AEV} \text{ ou } \ang{VEA}} & $E$ & $[EA)$ \ et \ $[EV)$ \\ \hline
e) \ \rep[3.3cm]{\ang{NRP} \text{ ou } \ang{PRN}} & $R$ & $[RN)$ \ et \ $[RP)$ \\ \hline
\end{tabular}
\end{minipage}
\end{exercicedevoir}

\vspace{-1.1cm}
\begin{exercicedevoir}[4][Entoure la bonne écriture]
\vspace{-0.3cm}
a) et b) : entoure \textbf{un côté} de l'angle. \quad c) et d) : entoure \textbf{les deux côtés}. \hfill \textit{\small (1 point par réponse)}

\renewcommand{\arraystretch}{1.25}
\begin{tabular}{@{}l@{\hspace{1.5cm}}l@{}}
\textit{Ex.} $\ang{GOK}$ : \ $[GO)$ \ \entoureex{$[OG)$} \ $[OG]$ \ $[KG)$ & \\
a) $\ang{LMN}$ : \ \entoure{$[LM)$} \ \entoure[r]{$[ML)$} \ \entoure{$[ML]$} \ \entoure{$[NM)$} &
b) $\ang{UVW}$ : \ \entoure{$[WV)$} \ \entoure{$[UW)$} \ \entoure{$[VW]$} \ \entoure[r]{$[VW)$} \\
\multicolumn{2}{@{}l@{}}{c) Les côtés de $\ang{CDE}$ : \quad \entoure{$[CD)$ et $[CE)$} \quad \entoure{$[DC]$ et $[DE]$} \quad \entoure[r]{$[DC)$ et $[DE)$}} \\
\multicolumn{2}{@{}l@{}}{d) Les côtés de $\ang{PSQ}$ : \quad \entoure[r]{$[SP)$ et $[SQ)$} \quad \entoure{$[PS)$ et $[PQ)$} \quad \entoure{$[SP)$ et $[QS)$}} \\
\end{tabular}
\end{exercicedevoir}

\vspace{-1.1cm}
\begin{exercicedevoir}[8][Sur une figure]
\vspace{-0.3cm}
Complète, puis repasse en couleur les côtés de l'angle demandé.
\hfill \textit{\small (a, b, d, e : 1,5 point ; c, f : 1 point)}

\begin{minipage}[c]{0.3\linewidth}\centering\eventail\end{minipage}\hfill
\begin{minipage}[c]{0.68\linewidth}
\renewcommand{\arraystretch}{1.35}
\begin{tabular}{@{}l@{}}
a) $\ang{DTF}$ : sommet \rep{T} ; \ côtés \rep[1.3cm]{[TD)} et \rep[1.3cm]{[TF)} \\
b) $\ang{FTG}$ : sommet \rep{T} ; \ côtés \rep[1.3cm]{[TF)} et \rep[1.3cm]{[TG)} \\
c) Repasse en couleur les côtés de l'angle $\ang{GTD}$. \\
\end{tabular}
\end{minipage}

\begin{minipage}[c]{0.3\linewidth}\centering\triangleCotes{G}{B}{K}{3,1.9}{0,0.4}{2.4,-0.3}{320}{95}{300}\end{minipage}\hfill
\begin{minipage}[c]{0.68\linewidth}
\renewcommand{\arraystretch}{1.35}
\begin{tabular}{@{}l@{}}
d) $\ang{GBK}$ : sommet \rep{B} ; \ côtés \rep[1.3cm]{[BG)} et \rep[1.3cm]{[BK)} \\
e) $\ang{BKG}$ : sommet \rep{K} ; \ côtés \rep[1.3cm]{[KB)} et \rep[1.3cm]{[KG)} \\
f) Repasse en couleur les côtés de l'angle $\ang{BGK}$. \\
\end{tabular}
\end{minipage}
\end{exercicedevoir}

\vspace{-0.6cm}
{\small\competences{
6C04-GEO11 & Sommet et côtés d'un angle & & & & \\ \hline
}}
}

\begin{document}

\corrigefalse
\feuille

\end{document}

%%% Local Variables:
%%% mode: LaTeX
%%% TeX-master: t
%%% End:
